\documentclass[12pt]{article}
\usepackage[a4paper,margin=2.5cm]{geometry}
\usepackage[T1]{fontenc}
\usepackage[utf8]{inputenc}
\usepackage{lmodern}
\usepackage{microtype}
\usepackage{xcolor}
\usepackage{listings}
\usepackage[hidelinks]{hyperref}
\definecolor{codegray}{RGB}{245,245,245}
\lstset{basicstyle=\ttfamily\small,backgroundcolor=\color{codegray},frame=single,
  breaklines=true,showstringspaces=false,columns=fullflexible,keepspaces=true}
\setlength{\parindent}{0pt}
\setlength{\parskip}{0.55em}

\begin{document}
\pagenumbering{gobble}
\begin{titlepage}
\centering
{\large Ministerul Educa\c{t}iei \c{s}i Cercet\u{a}rii al Republicii Moldova\par}
\vspace{0.45cm}
{\large Universitatea Tehnic\u{a} a Moldovei\par}
\vspace{0.45cm}
{\large Facultatea Calculatoare, Informatic\u{a} \c{s}i Microelectronic\u{a}\par}
\vspace{0.45cm}
{\large Departamentul Ingineria Software \c{s}i Automatic\u{a}\par}
\vspace{3cm}
{\Large\bfseries Report on Laboratory Work No. 8\par}
\vspace{0.55cm}
{\LARGE\bfseries Process Calls and Pipes Inside xv6\par}
\vspace{0.35cm}
{\large Operating Systems\par}
\vfill
\begin{flushleft}
\textbf{Done by:} st. gr. FAF-241, Ca\c{t} Chiril\\[0.8cm]
\textbf{Checked by:} asist. univ. Rei\c{t}man Patricia
\end{flushleft}
\vfill
{Chi\c{s}in\u{a}u -- 2026\par}
\end{titlepage}

\pagenumbering{roman}
\tableofcontents
\clearpage
\pagenumbering{arabic}

\section{Purpose and environment}
I ran process creation, program replacement, and two-way communication
inside xv6. I added \texttt{user/run.c} and \texttt{user/pingpong.c} to the
integrated \texttt{os/xv6-riscv} tree and its \texttt{UPROGS} list, then
built the image and booted it with QEMU on 26 October 2026. Matching code
copies are in \texttt{lab8/code/}. The required \texttt{notes.md} records
the Observe answers and command output. The optional \texttt{primes}
exercise was not part of the required deliverable.

\section{Fork, exec, and wait inside xv6}
The final \texttt{run.c} prints its PID, forks, executes \texttt{echo} in
the child, and waits for the child's status in the parent:

\begin{lstlisting}
$ run
parent: my pid is 3
child: my pid is 4
hello from the child
parent: child 4 ended with status 0
\end{lstlisting}

I also compiled and ran two temporary variations. With
\texttt{exec("ls", args)}, the child printed the xv6 directory listing
and exited successfully:

\begin{lstlisting}
$ run
parent: my pid is 3
child: my pid is 4
.              1 1 1024
..             1 1 1024
README         2 2 2441
... (other directory entries)
console        3 29 0
parent: child 4 ended with status 0
\end{lstlisting}

With \texttt{exec("nosuchprog", args)}, replacement failed, so the child
continued at the following error line:

\begin{lstlisting}
$ run
parent: my pid is 3
child: my pid is 4
exec failed
parent: child 4 ended with status 1
\end{lstlisting}

The xv6 program uses \texttt{fork}, \texttt{exec}, and \texttt{wait}; the
Lab 7 POSIX program uses \texttt{fork}, \texttt{execlp}, and
\texttt{waitpid}. Both sequences create a child, replace its program, and
collect its status. \texttt{echo}, rather than the original \texttt{run}
code, prints \texttt{hello from the child} after the successful
\texttt{exec}. When replacement fails, xv6's program calls
\texttt{exit(1)}, whereas the Lab 7 program calls \texttt{exit(127)}.
Those explicit choices explain the different parent-reported statuses.

\clearpage
\section{A two-pipe pingpong}
\texttt{pingpong.c} creates a parent-to-child pipe and a child-to-parent
pipe before forking. The parent sends one byte, the child reads it and
prints \texttt{received ping}, then writes it back. The parent reads the
reply and prints \texttt{received pong}:

\begin{lstlisting}
$ pingpong
6: received ping
5: received pong
$ pingpong
8: received ping
7: received pong
\end{lstlisting}

Two pipes are needed because each pipe is a one-way byte stream. One
carries the request to the child and the other carries the reply to the
parent. Each process closes the ends it does not use. The parent's
\texttt{read(c2p[0], \&b, 1)} blocks until the child writes the reply;
this call orders the two messages. The later \texttt{wait(\&status)}
collects the child's exit status.

\section{Conclusion}
The xv6 run followed the same process lifecycle observed with POSIX calls
in Lab 7: \texttt{fork} creates a child, \texttt{exec} runs a different
program in it, and \texttt{wait} reaps it. The failed \texttt{exec} returned
to the child and exposed its explicit exit code. The two-pipe experiment
showed how blocking reads coordinate a request and reply without a timing
delay.

\end{document}
